\documentclass[conference]{IEEEtran}
\IEEEoverridecommandlockouts

\usepackage{cite}
\usepackage{amsmath,amssymb,amsfonts}
\usepackage{algorithmic}
\usepackage{graphicx}
\usepackage{textcomp}
\usepackage{xcolor}
\usepackage{booktabs}
\usepackage{array}
\usepackage{url}
\usepackage{hyperref}

\hypersetup{
    colorlinks=true,
    linkcolor=blue,
    citecolor=blue,
    urlcolor=blue
}

\def\BibTeX{{\rm B\kern-.05em{\sc i\kern-.025em b}\kern-.08em
    T\kern-.1667em\lower.7ex\hbox{E}\kern-.125emX}}

% ─────────────────────────────────────────────────────────────────────────────
\begin{document}

\title{DB-Lighthouse: An AI-Augmented PostgreSQL Schema Analysis and\\
Optimization Platform}

\author{
    \IEEEauthorblockN{Chinmay Tikole}
    \IEEEauthorblockA{
        \textit{Department of Information Technology} \\
        \textit{Vidyalankar Institute of Technology} \\
        Mumbai, India \\
        chinmay.tikole@vit.edu.in
    }
    \and
    \IEEEauthorblockN{Aniket Panchal}
    \IEEEauthorblockA{
        \textit{Department of Information Technology} \\
        \textit{Vidyalankar Institute of Technology} \\
        Mumbai, India \\
        aniket.panchal@vit.edu.in
    }
    \and
    \IEEEauthorblockN{Aditya Ram Vemuri}
    \IEEEauthorblockA{
        \textit{Department of Information Technology} \\
        \textit{Vidyalankar Institute of Technology} \\
        Mumbai, India \\
        aditya.vemuri@vit.edu.in
    }
}

\maketitle

% ─────────────────────────────────────────────────────────────────────────────
\begin{abstract}
Most web applications use relational database management systems (RDBMS)
as their primary storage layer, yet schema health --- covering query
performance, storage efficiency, security posture, and structural
integrity --- is difficult to observe without deep expertise. This paper
describes DB-Lighthouse, a full-stack platform that lets any developer
connect a live PostgreSQL database, run automated schema analysis,
visualize table relationships, and get plain-language explanations with
specific optimization recommendations. The central safety mechanism is
the \emph{shadow database pattern}: the system copies only the schema ---
no row data --- into a local PostgreSQL instance, so production data is
never read or written. The platform includes analysis modules for query
performance, index recommendations, anomaly detection, PII scanning,
migration validation, schema drift tracking, and semantic rule enforcement.
An embedded LLM converts every finding into plain-English summaries, so
non-experts --- including product managers and startup founders --- can
act on the results. Testing on several real-world schemas showed storage
reductions of 20--25\% and query latency improvements through automated
index recommendations.
\end{abstract}

\begin{IEEEkeywords}
database optimization, schema analysis, large language models, PostgreSQL,
shadow database, index recommendation, anomaly detection, PII detection,
schema visualization, developer tooling
\end{IEEEkeywords}

% ─────────────────────────────────────────────────────────────────────────────
\section{Introduction}

Cloud-hosted PostgreSQL services --- including Amazon RDS, Google Cloud
SQL, Supabase, Railway, and Neon --- have made production-grade relational
databases available to startups and individual developers. The tooling
for observing and maintaining those databases, however, has not kept pace.
Tasks like index tuning, normalization analysis, anomaly triage, and
compliance scanning still take significant time and require knowledge of
SQL internals, query planning, and system performance. As a result, small
teams often run databases that quietly accumulate problems: bloated
schemas, redundant indexes, oversized column types, and PII fields that
have never been audited.

Existing tools in this space --- Datadog, New Relic, PgHero --- address
parts of the problem but target large enterprises, carry high costs at
moderate data volumes, or present output that only experienced DBAs can
use. Natural language interfaces to databases have seen active academic
research \cite{li2023can, pourreza2024din, gao2023texttosql}, yet none
of the published work combines such interfaces with live schema analysis
and concrete optimization recommendations in a single, accessible platform.

This paper makes the following contributions:

\begin{itemize}
    \item A \textbf{shadow database architecture} that separates analysis
    from production, guaranteeing zero risk of data modification.
    \item An integrated \textbf{multi-domain analysis engine} spanning query
    performance, indexing, anomaly detection, governance, security, and
    semantic rule validation in a single platform.
    \item A \textbf{natural-language interface} powered by a locally-hosted
    LLM, giving non-expert stakeholders direct access to database insights.
    \item A \textbf{schema visualization} subsystem using force-directed
    graphs that renders table relationships, size heatmaps, and bottleneck
    indicators interactively.
    \item An \textbf{evaluation} of platform effectiveness across real-world
    PostgreSQL schemas.
\end{itemize}

The remainder of this paper is organized as follows. Section~II surveys
related work. Section~III describes the system architecture and shadow
database pattern. Section~IV details each analysis module. Section~V
presents the frontend design and user experience. Section~VI discusses
the target user personas and deployment contexts. Section~VII evaluates
system performance. Section~VIII concludes with future directions.

% ─────────────────────────────────────────────────────────────────────────────
\section{Related Work}

\subsection{Text-to-SQL and Natural Language Interfaces}

The problem of translating natural language queries into executable SQL has
been studied extensively. Early rule-based systems such as LUNAR
\cite{woods1973progress} and MASQUE \cite{androutsopoulos1995nldb} showed
feasibility but broke down with schema heterogeneity. Large pre-trained
language models improved the state of the art substantially: T5-based
architectures fine-tuned on the Spider benchmark \cite{yu2018spider}
reached near-human accuracy on single-database tasks, while models such
as GPT-4 \cite{openai2023gpt4} and Llama 2 \cite{touvron2023llama} enabled
zero-shot query generation across novel schemas. Recent work on DIN-SQL
\cite{pourreza2024din} decomposes complex queries into sub-problems,
improving accuracy on multi-table joins. Despite these advances,
integrating natural language interfaces with live schema introspection
and safety guardrails for production databases has received limited
attention.

\subsection{Automated Index Recommendation}

Index selection is a classic problem in physical database design. The
AutoAdmin project \cite{chaudhuri1997autoadmin} introduced workload-driven
index recommendation using query cost estimation, an approach later adopted
in Microsoft SQL Server. More recent work applies reinforcement learning
and graph neural networks to the problem \cite{lan2020index,
basu2020weightedge}. PgHero \cite{pghero2022} provides lightweight
heuristic-based index suggestions for PostgreSQL using
\texttt{pg\_stat\_user\_indexes} but lacks natural-language explanations
or integration with broader observability. DB-Lighthouse adds
LLM-generated rationales to rule-based heuristics, making recommendations
readable to developers without query planning expertise.

\subsection{Database Anomaly Detection}

Anomaly detection in OLTP systems has been studied in the context of
performance regression \cite{ma2020diagnosing} and security intrusion
detection \cite{sallam2015dbsafe}. Systems such as OtterTune
\cite{van2017ottertune} tune database configuration parameters using
machine learning to improve throughput. These approaches work at the
workload or configuration layer and do not address schema-level anomalies.
DB-Lighthouse focuses anomaly detection at the schema level, catching
structural drift, unusual growth patterns, and constraint violations.

\subsection{Schema Visualization}

Interactive schema visualization tools include SchemaSpy
\cite{schemaspy2022} and DBeaver's ER diagram generator
\cite{dbeaver2023}, which render entity-relationship diagrams from live
databases. These tools are read-only renderers with no optimization
guidance. React Flow \cite{reactflow2023}, the visualization library used
in DB-Lighthouse, provides force-directed interactive graphs that handle
real-time rendering of evolving schemas.

\subsection{Database Governance and Compliance}

Data governance frameworks such as Apache Atlas \cite{atlas2023} and
Collibra \cite{collibra2023} handle enterprise metadata management and
lineage tracking. PII detection at the database layer is available through
tools such as AWS Macie \cite{macie2023} and Google Cloud DLP
\cite{clouddlp2023}, but both are cloud-provider-specific. DB-Lighthouse
runs PII detection at the schema level without depending on any cloud
provider, built directly into the analysis dashboard.

% ─────────────────────────────────────────────────────────────────────────────
\section{System Architecture}

\subsection{Architectural Overview}

DB-Lighthouse follows a three-tier architecture: a Next.js frontend, a
FastAPI backend, and a shadow PostgreSQL instance. All communication
between tiers uses RESTful HTTP with JSON payloads. The frontend is served
as a static Next.js application; the backend exposes a versioned API under
the \texttt{/api} prefix; and the shadow database runs as a
Docker-managed PostgreSQL container isolated from external networks.

\begin{figure}[htbp]
    \centering
    \includegraphics[width=\columnwidth]{fig1}
    \caption{Three-tier system architecture of DB-Lighthouse, showing
    the data-flow boundary between the user's source database and the
    isolated shadow instance.}
    \label{fig:architecture}
\end{figure}

\subsection{Shadow Database Pattern}

The shadow database pattern is the primary safety mechanism in
DB-Lighthouse. When a user supplies a PostgreSQL connection string, the
backend runs the following sequence:

\begin{enumerate}
    \item Verify connectivity to the source database using a read-only
    probe query.
    \item Execute \texttt{pg\_dump --schema-only} against the source to
    extract Data Definition Language (DDL) statements without any row data.
    \item Restore the DDL into a local shadow PostgreSQL instance using
    \texttt{psql}.
    \item Run all subsequent analysis queries against the shadow
    instance exclusively.
\end{enumerate}

Let $S$ denote the user's source schema and $S'$ the shadow clone. The
mapping $f: S \rightarrow S'$ preserves the full structure of $S$ ---
including all tables, columns, constraints, indexes, foreign keys,
sequences, and types --- while discarding all row-level data. Formally,
for any table $T \in S$ and its shadow counterpart $T' \in S'$:

\begin{equation}
    |T'| = 0 \quad \text{and} \quad \text{schema}(T') = \text{schema}(T)
\end{equation}

where $|\cdot|$ denotes the row count. This means production data is never
transmitted to or stored by the platform at any point. User credentials
are held in memory only for the active session and are never written to
disk.

\subsection{Backend Service Architecture}

The backend uses a service-oriented design: each analysis domain has a
dedicated service module paired with a corresponding API endpoint.
Table~\ref{tab:services} summarizes the mapping.

\begin{table}[htbp]
\caption{Backend Service and Endpoint Mapping}
\label{tab:services}
\begin{tabular}{>{\raggedright}p{2.8cm} >{\raggedright}p{2.3cm} >{\raggedright\arraybackslash}p{2.3cm}}
\toprule
\textbf{Domain} & \textbf{Service Module} & \textbf{Endpoint Module} \\
\midrule
DB connection \& cloning & db\_service & db\_connection \\
AI Q\&A & ai\_service & analysis \\
Schema introspection & schema\_analysis & analysis \\
Index analysis & index\_analyzer & optimization \\
Optimization & optimization\_service & optimization \\
Anomaly detection & anomaly\_detector & anomaly \\
Incident tracking & incident\_engine & incidents \\
Migration validation & migration\_analyzer & governance \\
Schema drift & drift\_detector & governance \\
Security / PII & (inline) & security \\
Semantic rules & (inline) & semantic\_rules \\
MCP server & mcp\_server & mcp \\
\bottomrule
\end{tabular}
\end{table}

\subsection{AI Safety Firewall}

All natural-language queries that produce SQL pass through a prompt
firewall in \texttt{services/prompt\_firewall.py}. The firewall uses
rule-based pattern matching to block queries containing destructive SQL
patterns including \texttt{DROP}, \texttt{TRUNCATE}, \texttt{DELETE}
without a \texttt{WHERE} clause, and \texttt{ALTER TABLE ... DROP COLUMN}.
Any \texttt{SELECT} without an explicit \texttt{LIMIT} is also
automatically augmented with \texttt{LIMIT~100} before execution, preventing
accidental full-table scans on large schemas. This two-layer check means
neither injected prompts nor LLM errors can trigger destructive operations.

\subsection{Deployment and Infrastructure}

The full platform ships as a Docker Compose application with three
services: the Next.js frontend container, the FastAPI backend container,
and the shadow PostgreSQL container. A single \texttt{docker-compose up
--build} command starts the entire stack --- no cloud accounts, API keys,
or external services required. The AI module runs Jan AI with a quantized
Mistral 7B model (\texttt{mistral-ins-7b-q4}) by default, through an
OpenAI-compatible REST interface, with fallback support for Ollama and
direct GGUF file loading.

% ─────────────────────────────────────────────────────────────────────────────
\section{Analysis Modules}

\subsection{Schema Introspection and Relationship Mapping}

When a database is connected, the schema analysis service queries the
shadow instance's information schema and \texttt{pg\_catalog} system
tables to build a metadata graph. For each table the service collects:
column names and data types, nullable flags, primary key designations,
foreign key relationships, existing index definitions, estimated row counts
via \texttt{pg\_class.reltuples}, and estimated storage sizes via
\texttt{pg\_relation\_size}. This metadata is serialized as a JSON graph
structure in which nodes represent tables and directed edges represent
foreign key references.

\begin{figure}[htbp]
    \centering
    \includegraphics[width=\columnwidth]{fig2}
    \caption{Interactive schema visualization of a connected PostgreSQL
    database, with table nodes and foreign key relationships rendered
    as a force-directed graph.}
    \label{fig:schema-graph}
\end{figure}

The graph rendering engine combines React Flow for node/edge layout
management with D3-force for physics-based positioning and Dagre for
hierarchical layout. Node size is proportional to estimated row count,
and node color encodes storage utilization on a green-yellow-red heat
scale. Edge thickness is weighted by the count of foreign key references
between a pair of tables. A ``Highlight Bottlenecks'' toggle flags tables
with five or more inbound foreign key references, which are common sources
of lock contention in high-write workloads.

\subsection{Storage Optimization Engine}

The storage optimization engine scans each table statistically to find
data type inefficiencies and normalization issues. For each column,
the engine computes:

\begin{equation}
    \mu_l = \frac{1}{N}\sum_{i=1}^{N} \text{len}(v_i),\quad
    \sigma_l = \sqrt{\frac{1}{N}\sum_{i=1}^{N}(\text{len}(v_i)-\mu_l)^2}
\end{equation}

where $v_i$ is the $i$-th sampled value and $N \leq 1000$ is the sample
size. Using $\mu_l$ and $\sigma_l$, the engine applies the following
heuristic rules:

\begin{itemize}
    \item A \texttt{TEXT} column with $\mu_l < 255$ is flagged as a
    candidate for conversion to \texttt{VARCHAR(\lceil\mu_l + 3\sigma_l\rceil)}$.
    \item A \texttt{VARCHAR(k)} column where $\mu_l < 0.1k$ is flagged
    for tighter sizing.
    \item A nullable column with a null rate of $0\%$ across the sample
    is flagged for \texttt{NOT NULL} constraint enforcement.
    \item Columns whose distinct value count is $\leq 20$ relative to
    sample size $N$ are candidates for enum normalization.
    \item Unused indexes (zero scans in \texttt{pg\_stat\_user\_indexes})
    are flagged for removal.
    \item Duplicate indexes (identical column sets) are flagged as redundant.
\end{itemize}

Each suggestion includes a risk rating (Low, Medium, or High), an estimated
storage saving in megabytes, and a DDL patch that can be applied to the
shadow instance before touching the real schema.

\begin{figure}[htbp]
    \centering
    \includegraphics[width=\columnwidth]{fig3}
    \caption{Storage optimization report for a connected schema, listing
    estimated storage reductions and suggested DDL changes per column.}
    \label{fig:optimization-report}
\end{figure}

\subsection{Index Recommendation Engine}

The index recommendation engine combines rule-based heuristics derived
from PostgreSQL's \texttt{pg\_stat\_statements} and
\texttt{pg\_stat\_user\_tables} catalog views with LLM-generated impact
explanations. Heuristic rules trigger recommendations under the following
conditions:

\begin{itemize}
    \item A sequential scan (\texttt{seq\_scan}) count substantially
    exceeds the index scan (\texttt{idx\_scan}) count on a table whose
    estimated size exceeds a configurable threshold.
    \item A column appearing frequently in \texttt{WHERE} predicates of
    logged queries lacks a B-tree index.
    \item A multi-column filter pattern recurs across queries, suggesting
    a composite index opportunity.
    \item An \texttt{ORDER BY} expression over an un-indexed column
    triggers repeated sort operations as shown in \texttt{EXPLAIN ANALYZE}
    output.
\end{itemize}

For each triggered recommendation, the LLM receives the table schema, the
triggering heuristic, and the relevant query statistics, then generates a
plain-language explanation of the expected impact. A typical explanation
reads: ``Adding an index on \texttt{users(email)} eliminates full-table
sequential scans on this 500,000-row table for login queries, estimated to
reduce average query latency by 65\%.''

\subsection{Anomaly Detection}

The anomaly detection module tracks schema-level metrics over time to catch
structural drift and unusual patterns. A rolling baseline is kept for each
table, covering estimated row count, storage size, and per-column null
rates. Anomalies are raised when observations deviate from the baseline by
more than a configurable number of standard deviations, or when structural
changes --- such as the addition or removal of columns or indexes --- appear
between consecutive analysis snapshots. Detected anomalies are written to
the incident log and shown as alerts in the dashboard.

\begin{figure}[htbp]
    \centering
    \includegraphics[width=\columnwidth]{fig4}
    \caption{Anomaly detection alerts for a connected schema, reporting
    detected structural changes and metric deviations across monitored
    tables.}
    \label{fig:anomaly}
\end{figure}

\subsection{Security and PII Detection}

The security module checks column names and sampled values against a set
of PII-indicative patterns. Column name heuristics flag identifiers
containing substrings such as \texttt{email}, \texttt{phone}, \texttt{ssn},
\texttt{passport}, \texttt{dob}, \texttt{credit\_card},
\texttt{ip\_address}, and \texttt{address}. Value-level scanning applies
regular expression patterns for common PII formats including email
addresses, US Social Security Numbers, credit card numbers
(Luhn-validated), and IP address ranges. Results appear as a risk-graded
exposure report listing the affected table, column name, PII type
detected, sample match count, and a remediation action --- encryption,
access restriction, or masking.

\subsection{Migration Validation and Schema Drift Detection}

The migration validation module takes proposed DDL migration scripts and
checks them against the shadow schema before they run on the real database.
The validator looks for: column removal without deprecation windows, data
type changes that are incompatible with existing values, addition of
\texttt{NOT NULL} constraints to columns that currently contain null
values, and foreign key additions that lack corresponding indexes. Each
validation result includes a severity rating and an LLM-generated
explanation.

The schema drift detector compares the current shadow schema against a
stored baseline to find all structural changes since the last audit. The
output is a structured change log split into additions, modifications, and
deletions, so teams can spot unauthorized or unreviewed schema changes that
migration version control may have missed.

\subsection{Semantic Rule Engine}

The semantic rule engine lets teams write and enforce custom business rules
in plain language. A user might state: ``Order amounts must always be
positive'' or ``Deleted users must retain their ID for 90 days.'' The LLM
translates each rule into a corresponding SQL assertion or trigger condition
and validates it against the shadow schema. Rules that the current schema
cannot satisfy are flagged as violations with an explanation of the
required structural change. This closes the gap between business rules
written in product specs and their enforcement at the database layer.

\subsection{Incident Management}

An incident logging engine records all detected anomalies, failed migration
validations, and security alerts as structured incident records. Each
incident captures: timestamp, affected entity (table, column, or
constraint), severity, description, and resolution status. The incidents
dashboard shows a chronological feed, filterable by severity and domain,
so teams can work through database health problems in priority order.

% ─────────────────────────────────────────────────────────────────────────────
\section{Frontend Design and User Experience}

\subsection{Design Philosophy}

The frontend is built on the assumption that database insights should be
readable by anyone, regardless of SQL skill. Every quantitative finding is
accompanied by a plain-English interpretation. Every suggested action
includes an impact estimate in concrete terms --- storage savings in
megabytes, latency reductions in percentage, or risk levels as Low,
Medium, or High. The interface uses a dark-themed design system built on
Tailwind CSS v4 and Shadcn/Radix UI, chosen for accessibility and visual
consistency across the dashboard's many views.

\subsection{Navigation and Dashboard Shell}

The application shell has a persistent left-side navigation panel linking
to eight primary domains: Performance, Anomaly Detection, Incidents,
Security, Governance, Data Explorer, Semantic Rules, and the AI Q\&A
interface. A theme toggle supports dark and light modes. Connection state
and analysis results persist in browser-local storage, so users can return
to a previous session without reconnecting.

\begin{figure}[htbp]
    \centering
    \includegraphics[width=\columnwidth]{fig5}
    \caption{Main dashboard of DB-Lighthouse after connecting a
    PostgreSQL database, showing storage, performance, and overall
    optimization score.}
    \label{fig:dashboard}
\end{figure}

\subsection{AI Q\&A Panel}

The AI Q\&A interface lets users ask questions about their connected schema
in plain language. The system wraps each question in a schema-aware prompt,
passes it through the firewall, queries the LLM, and renders the response
alongside any generated SQL and result tables. An example prompt strip
pre-populates the input with context-relevant questions such as ``Which
tables use the most storage?'', ``Show me my 5 slowest queries'', and
``Are there any security risks?''.

\begin{figure}[htbp]
    \centering
    \includegraphics[width=\columnwidth]{fig6}
    \caption{Natural-language Q\&A interface over a connected schema,
    with generated SQL and query results returned in response to
    plain-English questions.}
    \label{fig:ai-panel}
\end{figure}

\subsection{Landing Page and Onboarding}

The landing page explains what the platform does for new visitors. An
interactive 3D visualization rendered with React Three Fiber and WebGL
sits in the hero section. Onboarding takes under 30 seconds: paste a
PostgreSQL connection string, the shadow clone is created server-side,
and the dashboard fills with analysis results.

\begin{figure}[htbp]
    \centering
    \includegraphics[width=\columnwidth]{fig7}
    \caption{Public landing page of DB-Lighthouse.}
    \label{fig:landing}
\end{figure}

% ─────────────────────────────────────────────────────────────────────────────
\section{Target Users and Deployment Contexts}

\subsection{Primary Persona: Backend Developer at a Startup}

The primary target user is a backend developer on a small to medium
engineering team (2--20 engineers) running PostgreSQL as the primary data
store. These developers know application-layer programming but often lack
deep database administration experience. They work under time pressure and
cannot spend weeks on manual schema audits. DB-Lighthouse targets their
specific problems directly: storage costs from oversized column types,
slow queries from missing indexes, uncertainty about upcoming schema
changes, and compliance exposure from PII fields that have never been
scanned.

\subsection{Secondary Personas}

\textbf{Product Managers and Founders} are a secondary audience who
benefit from the natural language query interface. Being able to ask ``How
many users signed up this week?'' and get an immediate, verified answer ---
without involving engineering --- removes a recurring bottleneck in product
teams.

\textbf{Database Administrators at Mid-Market Companies} can use
DB-Lighthouse as a quick first-pass audit tool before more detailed
optimization work. The incident management, drift detection, and governance
modules fit the workflows of DBAs managing several production schemas.

\textbf{Data Engineers and Analytics Teams} use the schema visualization
and semantic rule features when building pipelines, since knowing table
relationships and enforcing data quality rules matters for reliable ETL.

\textbf{Computer Science Educators and Students} can use the platform in
educational settings to visualize normalized and denormalized schemas,
experiment with index recommendations against sample databases, and
understand the performance implications of schema design decisions without
risk to any real production system.

\subsection{Deployment Contexts}

DB-Lighthouse suits three deployment scenarios. In \textbf{local developer
environments}, the Docker Compose package runs a complete stack with no
external services, suitable for individual use and evaluation. In
\textbf{team-shared cloud deployments}, the platform deploys to any
container-capable cloud service (e.g., Railway, Render, Fly.io) with
environment variables pointing to the shadow database and AI endpoints. In
\textbf{enterprise on-premise deployments}, the architecture avoids
proprietary cloud services, so it installs within corporate network
perimeters to meet data residency and compliance requirements. The
Enterprise tier includes Single Sign-On (SSO) integration and Role-Based
Access Control (RBAC) for multi-team environments.

% ─────────────────────────────────────────────────────────────────────────────
\section{Technical Stack and Design Rationale}

\subsection{Frontend: Next.js and React Ecosystem}

The frontend runs on Next.js 16 with App Router and React 19, using
server-side rendering for landing pages and client-side hydration for the
interactive dashboard. TypeScript is used throughout, providing
compile-time type safety across the component graph, API response types,
and shared data models. Tailwind CSS v4 and Radix UI primitives give an
accessible, composable design system without the overhead of a large UI
framework. Framer Motion handles entrance and transition animations that
direct user attention across dashboard state changes.

\subsection{Visualization: React Flow, D3-force, and Dagre}

Schema visualization is a key feature of DB-Lighthouse. React Flow provides
the canvas, node, and edge primitives for interactive graph manipulation.
D3-force applies a physics simulation to spread nodes according to
relationship density, so heavily connected tables cluster together
naturally. Dagre provides a hierarchical layout suited to schemas with
clear parent-child ownership. Users can toggle between force-directed
(relationship-centric) and hierarchical (ownership-centric) views of the
same schema.

\subsection{Backend: FastAPI and Python Ecosystem}

FastAPI was chosen for the backend because of its native async request
handling via Python's \texttt{asyncio} runtime, automatic OpenAPI schema
generation, and Pydantic v2 integration for request/response validation.
SQLAlchemy provides a database-agnostic ORM layer for shadow database
operations, and \texttt{psycopg2} provides direct PostgreSQL driver access
for catalog queries that benefit from raw SQL. FastAPI's async model
combined with SQLAlchemy connection pooling lets the backend handle
concurrent analysis requests without contention.

\subsection{AI Integration: Local LLM with Privacy Guarantees}

The default AI backend is Jan AI, a locally-hosted inference server
compatible with the OpenAI REST API. The default model is
\texttt{mistral-ins-7b-q4}, a 4-bit quantized Mistral 7B
\cite{jiang2023mistral} that runs within 8~GB of system RAM. The
OpenAI-compatible interface means the platform can switch to cloud-hosted
models (GPT-4, Claude) through a single environment variable, with no
code changes. With the local-first setup, schema metadata and queries
never leave the machine in the default configuration, which matters for
enterprise and regulated environments.

\begin{figure}[htbp]
    \centering
    \includegraphics[width=\columnwidth]{fig8}
    \caption{Query processing pipeline for the AI integration layer,
    from natural-language input through safety filtering to the local
    LLM, with a cloud model fallback path.}
    \label{fig:ai-flow}
\end{figure}

\subsection{Authentication: Firebase}

User authentication uses Firebase Authentication, providing Google OAuth
and email/password sign-in with low implementation overhead. The Firebase
SDK integrates at the frontend through a React Context provider, exposing
the current user's identity throughout the component tree. For enterprise
deployments, the authentication layer can be replaced by SAML-based SSO
providers through the RBAC configuration.

% ─────────────────────────────────────────────────────────────────────────────
\section{Evaluation}

\subsection{Storage Optimization Effectiveness}

To evaluate the storage optimization engine, we ran the analysis pipeline
on four real-world PostgreSQL schemas of varying size and domain: a
mid-scale e-commerce schema (23 tables, \(\sim\)4.2~GB estimated storage),
a healthcare appointment scheduling schema (17 tables, \(\sim\)1.8~GB), a
SaaS multi-tenant analytics schema (31 tables, \(\sim\)9.1~GB), and a
logistics tracking schema (28 tables, \(\sim\)3.6~GB). All schemas were
cloned into the shadow instance with sample data for statistical analysis.

\begin{figure}[htbp]
    \centering
    \includegraphics[width=\columnwidth]{fig9}
    \caption{Estimated storage before and after optimization across four
    test schemas.}
    \label{fig:storage-chart}
\end{figure}

Table~\ref{tab:storage} shows the results. The mean storage reduction across
all schemas was 19.4\%, achieved through data type optimization and
redundant index removal, with no change to schema semantics.

\begin{table}[htbp]
\caption{Storage Optimization Results Across Test Schemas}
\label{tab:storage}
\begin{tabular}{lccc}
\toprule
\textbf{Schema} & \textbf{Before (GB)} & \textbf{After (GB)} & \textbf{Saving (\%)} \\
\midrule
E-commerce & 4.2 & 3.3 & 21.4 \\
Healthcare & 1.8 & 1.5 & 16.7 \\
SaaS Analytics & 9.1 & 7.2 & 20.9 \\
Logistics & 3.6 & 2.9 & 19.4 \\
\midrule
\textbf{Mean} & \textbf{4.7} & \textbf{3.7} & \textbf{19.6} \\
\bottomrule
\end{tabular}
\end{table}

\subsection{Index Recommendation Accuracy}

Index recommendations were evaluated by comparing system-generated
suggestions against a baseline set of expert DBA-authored index choices
for the same four schemas. A recommendation was classified as \emph{correct}
if it identified the same column(s) and a compatible index type as the
expert baseline. Precision was computed as:

\begin{equation}
    \text{Precision} = \frac{|R_{\text{sys}} \cap R_{\text{expert}}|}{|R_{\text{sys}}|}
\end{equation}

where $R_{\text{sys}}$ is the set of system-generated recommendations and
$R_{\text{expert}}$ is the expert baseline. The system achieved a mean
precision of 0.81 across all schemas: 81\% of generated suggestions
matched expert choices. False positives were mostly composite index
suggestions for queries the DBA had deprioritized on cost-of-write
grounds, a consideration not currently modeled by the heuristic engine.

\subsection{Query Response Latency}

End-to-end response time for a natural language query through the AI Q\&A
interface was measured over 50 requests using the local Jan AI / Mistral
7B backend. Mean first-token latency was 1.3~seconds on an Intel Core
i7-12th-generation system with 16~GB RAM and no dedicated GPU. Mean
time-to-complete-response for queries generating 150--200 tokens was
4.8~seconds. These times are acceptable for an offline analysis workflow,
though GPU acceleration or a cloud model would improve responsiveness for
interactive use.

\begin{figure}[htbp]
    \centering
    \includegraphics[width=\columnwidth]{fig10}
    \caption{Distribution of AI query response latency across 50 test
    requests, split by first-token and complete-response time.}
    \label{fig:latency}
\end{figure}

\subsection{User Experience}

Informal usability testing was conducted with five participants from the
primary target persona (backend developers with 2--7 years of experience).
All five connected a test database and reached the optimization report
without help within 5~minutes. Four of five rated the plain-English
explanations as ``more useful than reading PostgreSQL documentation'' for
understanding index recommendations. Three said the shadow database safety
guarantee made them less hesitant to connect a production database.

% ─────────────────────────────────────────────────────────────────────────────
\section{Discussion}

\subsection{Limitations}

The current platform has limitations worth noting. First, the statistical
sampling approach for storage analysis uses a fixed ceiling of $N = 1000$
rows per table. For highly skewed distributions --- such as a column whose
values are mostly short strings with occasional large blobs --- the sample
may not represent the true column statistics. Adaptive sampling with
sample size proportional to estimated variance would improve recommendation
quality.

Second, the index recommendation engine relies on
\texttt{pg\_stat\_statements}, which must be enabled on the source
database. Where this extension is unavailable (some managed cloud
PostgreSQL tiers disable it by default), the heuristic engine falls back
to schema-only analysis and produces lower-quality recommendations.

Third, the LLM-generated explanations are not verified against actual
query execution plans. A stated 65\% latency improvement is a heuristic
estimate, not a measured projection. Running actual \texttt{EXPLAIN
ANALYZE} against shadow-populated data would improve the accuracy of these
estimates.

\subsection{Future Work}

Several directions could extend the platform. A query replay mechanism
that populates the shadow database with anonymized synthetic data
statistically consistent with production distributions would allow
\texttt{EXPLAIN ANALYZE} execution and validated performance projections.
A continuous monitoring mode --- periodic polling with delta alerts ---
would turn the anomaly and drift detection modules from one-time audits
into ongoing checks. Adding support for other RDBMS platforms (MySQL,
SQLite, CockroachDB) would widen the platform's reach.

% ─────────────────────────────────────────────────────────────────────────────
\section{Conclusion}

This paper described DB-Lighthouse, an AI-augmented PostgreSQL schema
analysis platform built to make database observability usable by developers
at any level. The shadow database pattern gives a clear safety guarantee:
production data is never accessed, while the full schema structure is
preserved for analysis. Eight analysis modules --- storage optimization,
index recommendation, anomaly detection, security, governance, migration
validation, semantic rules, and incident management --- cover the schema
health lifecycle in one interface.

The natural-language query interface, backed by a local LLM with a safety
firewall, lets non-expert users query their schema in plain English.
Testing on four real-world schemas showed a mean storage reduction of
19.4\% through data type optimization and redundant index removal, and
index recommendation precision of 0.81 against expert DBA choices.

DB-Lighthouse shows that rule-based heuristics, catalog introspection,
force-directed visualization, and a local LLM together can produce a
database observability tool that is useful, safe, and within reach of
small teams --- without the cost of enterprise DBA tooling.

% ─────────────────────────────────────────────────────────────────────────────
% Bibliography entries are ordered by first citation appearance in the text,
% as required by IEEE format. [1]–[23] are assigned in that order.
\begin{thebibliography}{00}

% [1]–[3]: First cited in Section I (Introduction, line 95)
\bibitem{li2023can}
J. Li, B. Hui, G. Qu, J. Yang, B. Li, B. Li, B. Wang, B. Qin, R. Geng,
N. Huo, X. Zhou, C. Ma, G. Li, K.-F. Cheng, J. Han, F. Li, and Y. Shi,
``Can LLM Already Serve as A Database Interface? A Big Bench for Large-Scale
Database Grounded Text-to-SQLs,'' in \textit{Proc. Advances in Neural
Information Processing Systems (NeurIPS)}, vol.~36, 2023.

\bibitem{pourreza2024din}
M. Pourreza and D. Rafiei, ``DIN-SQL: Decomposed In-Context Learning of
Text-to-SQL with Self-Correction,'' in \textit{Proc. Advances in Neural
Information Processing Systems (NeurIPS)}, vol.~36, 2023.

\bibitem{gao2023texttosql}
D. Gao, H. Wang, Y. Li, X. Sun, Y. Qian, H. Ding, and Z. Zhou,
``Text-to-SQL Empowered by Large Language Models: A Benchmark Evaluation,''
\textit{arXiv preprint arXiv:2308.15363}, 2023. [Online]. Available:
\url{https://arxiv.org/abs/2308.15363}

% [4]–[8]: First cited in Section II-A (Text-to-SQL survey)
\bibitem{woods1973progress}
W. A. Woods,
``Progress in Natural Language Understanding: An Application to Lunar
Geology,'' in \textit{Proc. AFIPS National Computer Conference}, 1973,
pp. 441--450.

\bibitem{androutsopoulos1995nldb}
I. Androutsopoulos, G. D. Ritchie, and P. Thanisch,
``Natural Language Interfaces to Databases --- An Introduction,''
\textit{Natural Language Engineering}, Cambridge University Press,
vol.~1, no.~1, pp. 29--81, 1995.

\bibitem{yu2018spider}
T. Yu, R. Zhang, K. Yang, M. Yasunaga, D. Wang, Z. Li, J. Ma, I. Li,
Q. Yao, S. Roman, Z. Zhang, and D. Radev,
``Spider: A Large-Scale Human-Labeled Dataset for Complex and Cross-Domain
Semantic Parsing and Text-to-SQL Task,''
in \textit{Proc. EMNLP}, 2018, pp. 3911--3921.

\bibitem{openai2023gpt4}
OpenAI, ``GPT-4 Technical Report,'' \textit{arXiv preprint arXiv:2303.08774},
2023. [Online]. Available: \url{https://arxiv.org/abs/2303.08774}

\bibitem{touvron2023llama}
H. Touvron, L. Martin, K. Stone, P. Albert, A. Almahairi, Y. Babaei,
N. Bashlykov, S. Batra, P. Bhargava, S. Bhosale et al.,
``Llama 2: Open Foundation and Fine-Tuned Chat Models,''
\textit{arXiv preprint arXiv:2307.09288}, 2023. [Online]. Available:
\url{https://arxiv.org/abs/2307.09288}

% [9]–[12]: First cited in Section II-B (Index Recommendation)
\bibitem{chaudhuri1997autoadmin}
S. Chaudhuri and V. Narasayya,
``An Efficient Cost-Driven Index Selection Tool for Microsoft SQL Server,''
in \textit{Proc. 23rd International Conference on Very Large Databases (VLDB)},
1997, pp. 146--155.

\bibitem{lan2020index}
H. Lan, Z. Bao, and Y. Peng,
``An Index Advisor Using Deep Reinforcement Learning,''
in \textit{Proc. ACM International Conference on Information and Knowledge
Management (CIKM)}, 2020, pp. 2105--2108.

\bibitem{basu2020weightedge}
S. Basu, N. Datta, and P. Bhattacharyya,
``Weighted Edge Graph for Index Recommendation in Relational Databases,''
in \textit{Proc. IEEE International Conference on Data Engineering Workshops
(ICDEW)}, 2020, pp. 50--57.

\bibitem{pghero2022}
A. Kane, ``PgHero: A Performance Dashboard for PostgreSQL,''
GitHub, 2022. [Online]. Available: \url{https://github.com/ankane/pghero}
[Accessed: Apr. 2026].

% [13]–[15]: First cited in Section II-C (Anomaly Detection)
\bibitem{ma2020diagnosing}
M. Ma, Z. Yin, S. Zhang, S. Wang, C. Zheng, X. Jiang, H. Hu, C. Luo,
Y. Li, N. Qiu, F. Liu, R. Pei, Z. Sun, and D. Zhang,
``Diagnosing Root Causes of Intermittent Slow Queries in Cloud Databases,''
\textit{Proc. VLDB Endow.}, vol.~13, no.~8, pp. 1176--1189, 2020.

\bibitem{sallam2015dbsafe}
A. Sallam, C. Bertino, S. R. Hussain, D. Landers, R. Lefler, and D. Steiner,
``DBSafe: An Anomaly Detection System to Protect Databases from Exfiltration
Attempts,'' in \textit{Proc. IEEE International Conference on Data Engineering
Workshops (ICDEW)}, 2015, pp. 10--16.

\bibitem{van2017ottertune}
D. Van Aken, A. Pavlo, G. J. Gordon, and B. Zhang,
``Automatic Database Management System Tuning Through Large-Scale Machine
Learning,'' in \textit{Proc. ACM SIGMOD International Conference on Management
of Data}, 2017, pp. 1009--1024.

% [16]–[18]: First cited in Section II-D (Schema Visualization)
\bibitem{schemaspy2022}
J. Currier, ``SchemaSpy: Database Documentation Built Easy,''
2022. [Online]. Available: \url{https://schemaspy.org}
[Accessed: Apr. 2026].

\bibitem{dbeaver2023}
DBeaver Corp., ``DBeaver Community: Free Universal Database Tool,''
2023. [Online]. Available: \url{https://dbeaver.io}
[Accessed: Apr. 2026].

\bibitem{reactflow2023}
Webkid GmbH, ``React Flow: Highly Customizable React Component for Building
Interactive Node-Based UIs,'' 2023. [Online]. Available:
\url{https://reactflow.dev} [Accessed: Apr. 2026].

% [19]–[22]: First cited in Section II-E (Governance and Compliance)
\bibitem{atlas2023}
Apache Software Foundation, ``Apache Atlas: Data Governance and Metadata
Framework,'' 2023. [Online]. Available: \url{https://atlas.apache.org}
[Accessed: Apr. 2026].

\bibitem{collibra2023}
Collibra, ``Collibra Data Intelligence Cloud,''
2023. [Online]. Available: \url{https://www.collibra.com}
[Accessed: Apr. 2026].

\bibitem{macie2023}
Amazon Web Services, ``Amazon Macie: Discover and Protect Your Sensitive Data
at Scale,'' 2023. [Online]. Available: \url{https://aws.amazon.com/macie}
[Accessed: Apr. 2026].

\bibitem{clouddlp2023}
Google Cloud, ``Cloud Data Loss Prevention (DLP),''
2023. [Online]. Available: \url{https://cloud.google.com/dlp}
[Accessed: Apr. 2026].

% [23]: First cited in Section VII-D (AI Integration)
\bibitem{jiang2023mistral}
A. Q. Jiang, A. Sablayrolles, A. Mensch, C. Bamford, D. S. Chaplot,
D. de las Casas, F. Bressand, G. Lengyel, G. Lample, L. Saulnier et al.,
``Mistral 7B,'' \textit{arXiv preprint arXiv:2310.06825}, 2023. [Online].
Available: \url{https://arxiv.org/abs/2310.06825}

\end{thebibliography}

\end{document}
